
\section{Potential Applications}\label{sec:app}
With the widespread utilization of networks as a data modeling structure for representing diverse relational information across social, natural, and academic domains, graph prompting has begun to move from general-purpose adaptation techniques toward concrete application scenarios. In this section, we primarily focus on applications where prompts are instantiated over graph-structured data or graph representations, while also noting related settings where graph knowledge supports prompt-based reasoning. We discuss its emerging applications in online social networks, recommender systems, knowledge management, and biology, with an emphasis on dynamic interaction graphs, user-item graphs, knowledge graphs, and biomedical graphs.


\textbf{Online Social Networks. }
Online social platforms consist of users who can be represented as nodes, and their social connections, reposting traces, or interaction events form online social networks (OSNs). Previous research has investigated the potential of prompting mechanisms in identifying fake news within OSNs to prevent malicious attacks. For example, \citet{wu2023promptandalign} apply textual prompts to pre-trained language models (PLMs) to distill general semantic information, which is then combined with propagation dynamics in social networks for improved classification. Beyond text-oriented prompting, OSNs can also be viewed as dynamic interaction graphs whose node states and edges evolve over time. \citet{chen2024prompt} show that prompt learning can be used to adapt temporal interaction representations to downstream tasks, and DyGPrompt \cite{yu2025dygprompt} further designs node-time conditional prompts to bridge both task-objective gaps and temporal-distribution gaps between dynamic graph pre-training and downstream inference. These advances make OSN applications less speculative: few-shot fake news detection, malicious-account identification, and interaction-level risk assessment can all benefit from prompts that encode temporal and structural context when labeled data are scarce and expensive to obtain.



\textbf{Recommender Systems. }
E-commerce platforms provide a valuable opportunity to leverage recommender systems for enhancing online services, because user-item interactions naturally form bipartite or heterogeneous graphs. Early graph prompt tuning studies have explored cross-domain recommendation and contrastive pre-training adaptation \cite{wu2023personalized, hao2024motifbased}. For cross-domain recommendation, \citet{yi2023contrastive} introduce extra prompt nodes when transferring a pre-trained recommendation model to a target domain, thereby improving domain adaptation with limited target supervision. Meanwhile, \citet{yang2023empirical} propose personalized user prompts to bridge the gap between contrastive pretext tasks \cite{wu2021selfsupervised, yu2022are} and downstream recommendation objectives. Recent work further extends this line from static adaptation to dynamic and streaming user-item graphs. GraphPro \cite{yang2024graphpro} combines graph pre-training with temporal and structural prompts to improve dynamic recommendation, while GPT4Rec \cite{zhang2024gpt4rec} introduces node-level, structure-level, and view-level prompts to guide streaming recommendation over evolving interaction graphs. These studies suggest that graph prompting is particularly suitable for recommender systems where user interests, item popularity, and interaction patterns change continuously, and where efficient adaptation is required for cold-start, cross-domain, and continual recommendation settings.



\textbf{Knowledge Management. }Knowledge graphs (KGs) provide another application scenario in which prompt tuning can be applied directly to graph-structured knowledge. A central line of work performs prompting on KGs using pre-trained KG models to facilitate knowledge transfer across different KG data and tasks. In this line, \citet{zhang2023structure} design structure-aware pre-training objectives and then employ prompt tuning to transfer task-relevant knowledge to downstream KG tasks. Recent KG completion studies further suggest that adaptive or dual prompting can help align structured relational patterns with task-specific prediction objectives: APKGC \cite{maoliniyazi2026apkgc} uses adaptive prompting for KG completion, and FFCL-KGC \cite{zhu2026ffclkgc} incorporates dual prompting into feature-fusion contrastive learning. In parallel, a related but broader direction combines grounded knowledge from KGs with the reasoning and language-generation abilities of LLMs \cite{wang2023knowledge, robinson2023leveraging, park2023graphguided, zhang2023benchmarking, pan2023unifying}. For example, \citet{tian2023graph} introduce graph neural prompting to adapt KG knowledge to LLMs in a time- and parameter-efficient manner. These studies indicate that graph prompting can support structured knowledge transfer within KGs, while KG-grounded prompting for LLMs provides a neighboring direction for graph-enhanced knowledge reasoning.





\textbf{Biology. }Molecules can be represented naturally as graphs, where atoms serve as nodes and chemical bonds act as edges \cite{rong2020selfsupervised}. Such graph modeling provides a basis for applying graph representation learning methods to tasks such as molecular property prediction, drug interaction analysis, and protein modeling \cite{rong2020selfsupervised, liu2023molca, zhao2023gimlet, edwards2023synergpt}. Earlier molecular graph learning methods often trained task-specific models on individual datasets \cite{rong2020selfsupervised, sun2020infograph}, which limited cross-task generalization. Although LLMs and language models have been explored as general-purpose tools for understanding molecules \cite{qian2023can, liu2023one}, text-only descriptions may overlook the internal graph structures of molecules \cite{zhang2023large}. Recent graph-language models attempt to preserve both structural dependencies and semantic knowledge under few-shot or zero-shot settings \cite{zhao2023gimlet, edwards2023synergpt}. For instance, GIT-Mol \cite{liu2023gitmol} uses LoRA \cite{hu2021lora} for efficient downstream adaptation in molecular LLM systems. Graph prompting now provides more direct biomedical evidence. MOAT \cite{long2024moat} designs atom-, geometry-, and task-level prompts for 3D molecular graphs; DDIPrompt \cite{wang2024ddiprompt} applies graph prompt learning to rare and imbalanced drug-drug interaction event prediction; and Prompt-DDG \cite{wu2024promptddg} uses microenvironment-aware hierarchical prompts to predict mutation effects in protein-protein interactions. Together with protein multimer prompt learning \cite{gao2024protein}, these studies show that graph prompting can support efficient task adaptation across molecular, drug-safety, and protein-interaction applications.

